\documentclass[11pt]{article}
\usepackage[margin=1in]{geometry}
\usepackage{amsmath,amssymb,amsthm,mathtools,hyperref}
\newtheorem{theorem}{Theorem}[section]
\newtheorem{corollary}[theorem]{Corollary}
\newtheorem{remark}[theorem]{Remark}
\DeclareMathOperator{\CH}{CH}
\title{The elementary-symmetric candidate on a Fermat fourfold\\is an ambient complete intersection}
\author{}
\date{July 2026}
\begin{document}
\maketitle

\begin{abstract}
Let $X_{3d}\subset\mathbf P^5_{\mathbf C}$ be the Fermat fourfold and let
$W_d$ be the surface proposed in Question~1 of da Silva's 2021 notes, defined
by the first three elementary symmetric polynomials in
$x_0^d,\ldots,x_5^d$.  Newton's identity proves scheme-theoretically that
$W_d$ is the intersection on $X_{3d}$ of two ambient divisors.  Consequently
\[
 [W_d]=2d^2H^2\in\CH^2(X_{3d}),
\]
so its primitive cohomology class and every primitive Fermat-character
projection vanish.  This gives a negative answer for the proposed candidate,
without deciding algebraicity of the exceptional characters.  We also record
the independent fact that any cycle class invariant under coordinatewise
$d$-th roots of unity can occur only in character sectors whose coordinates
are all divisible by $d$.
\end{abstract}

\section{The candidate}
Let $k$ be a field in which $3$ is invertible, let $d\geq1$, and put
\[
 S=k[x_0,\ldots,x_5],\qquad y_i=x_i^d,
 \qquad F_{3d}=\sum_{i=0}^5x_i^{3d}.
\]
Write $e_j=e_j(y_0,\ldots,y_5)$ for the elementary symmetric polynomials.
Define
\[
 X_{3d}=V(F_{3d})\subset\mathbf P^5_k,
 \qquad W_d=V(e_1,e_2,e_3)\subset\mathbf P^5_k.
\]
This is the candidate denoted $W$ in Question~1 of \cite{daSilva}.  The
condition $3\nmid d$ used there is not needed below.

\section{Scheme structure and cycle class}
\begin{theorem}\label{thm:ci}
If $3$ is invertible in $k$, then
\[
 (e_1,e_2,e_3)=(e_1,e_2,F_{3d})\subset S.
\]
This ideal is a complete-intersection ideal of height three.  Hence
\[
 W_d=V_{X_{3d}}(e_1,e_2)
\]
scheme-theoretically, and $W_d$ is a pure two-dimensional complete
intersection of type $(d,2d,3d)$ in $\mathbf P^5$.
\end{theorem}
\begin{proof}
Newton's third identity is
\[
 \sum_{i=0}^5y_i^3=e_1^3-3e_1e_2+3e_3.
\]
After $y_i=x_i^d$, this becomes
\begin{equation}\label{eq:newton}
 F_{3d}=e_1^3-3e_1e_2+3e_3.
\end{equation}
It gives both inclusions between the displayed ideals, the reverse inclusion
using division by $3$.

For completeness, let $A=k[y_0,\ldots,y_5]$.  The elementary symmetric
functions are algebraically independent, and $A$ is integral over
$k[e_1,\ldots,e_6]$.  Thus $(e_1,e_2,e_3)$ has height three in the
Cohen--Macaulay ring $A$, and its generators form a regular sequence.  The map
\[
 A\longrightarrow S,\qquad y_i\longmapsto x_i^d,
\]
is finite free, with basis $\prod_i x_i^{r_i}$ for $0\leq r_i<d$.
Flat base change preserves the regular sequence.  Equation
\eqref{eq:newton} then permits replacing $e_3$ by $F_{3d}$.  The claimed
height, purity, and scheme-theoretic description follow.
\end{proof}

Now take $k=\mathbf C$ and let
$H=c_1(\mathcal O_{X_{3d}}(1))$.
\begin{corollary}\label{cor:class}
In the Chow group of the smooth Fermat fourfold,
\[
 [W_d]=2d^2H^2\in\CH^2(X_{3d}),\qquad \deg W_d=6d^3.
\]
Its cohomology class has zero primitive component and therefore zero
projection to every primitive Fermat-character eigenspace.
\end{corollary}
\begin{proof}
By Theorem~\ref{thm:ci}, $W_d$ is the proper scheme-theoretic intersection on
$X_{3d}$ of Cartier divisors of classes $dH$ and $2dH$.  Refined
intersection, including scheme multiplicities \cite{Fulton}, gives
\[
 [W_d]=(dH)(2dH)=2d^2H^2.
\]
Since $\int_{X_{3d}}H^4=3d$, its degree is $6d^3$.
The cycle class lies on the ambient Lefschetz line $\mathbf QH^2$, whose
projection to $H^4_{\mathrm{prim}}(X_{3d},\mathbf Q)$ is zero.  Primitive
Fermat-character spaces lie in this primitive summand \cite{Shioda}.
\end{proof}

Reducedness and irreducibility of $W_d$ are neither asserted nor needed.
The conclusion concerns the primitive component of its cohomology class; the
total class $2d^2H^2$ is nonzero.

\section{A coordinate-power sector obstruction}
Let
\[
 X_m=\left\{\sum_{i=0}^N x_i^m=0\right\}\subset\mathbf P^N_{\mathbf C},
 \qquad d\mid m,
\]
and let
\[
 G_m=\mu_m^{N+1}/\mu_m^{\mathrm{diag}},\qquad
 K_d=\mu_d^{N+1}/\mu_d^{\mathrm{diag}}\subset G_m,
\]
where the inclusion is induced by $\mu_d\subset\mu_m$.
A character of $G_m$ is represented by
$\alpha=(a_0,\ldots,a_N)\in(\mathbf Z/m)^{N+1}$ with
$\sum_i a_i=0$.

\begin{theorem}\label{thm:sector}
If a cohomology class $v$ on $X_m$ is invariant under $K_d$, then its
$\alpha$-character projection is zero unless
\[
 d\mid a_i\quad\text{for every }i.
\]
In particular this applies to the fundamental class of a cycle defined by
homogeneous equations in $x_0^d,\ldots,x_N^d$.
\end{theorem}
\begin{proof}
Finite abelian group representations over $\mathbf C$ are semisimple.  If
$\pi_\alpha$ denotes the character projector, averaging a $K_d$-invariant
class first over $K_d$ shows
\[
 \pi_\alpha(v)=0
 \quad\text{unless}\quad \chi_\alpha|_{K_d}=1.
\]
Explicitly,
\[
 \chi_\alpha|_{K_d}[(\zeta_i)]=\prod_i\zeta_i^{a_i}.
\]
This is well defined on the diagonal quotient because $\sum_i a_i=0\pmod m$.
Setting all but one coordinate equal to $1$ shows that the restriction is
trivial exactly when every $a_i$ is divisible by $d$.
Equations in the $x_i^d$ are fixed by $K_d$, so their scheme and fundamental
cycle are invariant.
\end{proof}

For the fourfold $X_{3d}$, a primitive $(2,2)$ character surviving this test
must be a permutation of
\[
 (d,d,d,2d,2d,2d).
\]
Indeed, primitive Fermat characters have all six coordinates nonzero, and the
$(2,2)$ condition is
\[
 \sum_i\langle a_i\rangle=3(3d)=9d
\]
for least-positive representatives \cite{Shioda}.  Divisibility by $d$ makes
each coordinate $d$ or $2d$.  If $r$ coordinates are $2d$, their sum is
$(6+r)d$, so $r=3$.

\section{Degree 33 and the scope of the answer}
For $d=11$ the complete-intersection computation gives
\[
 [W_{11}]=242H^2,\qquad \deg W_{11}=7986.
\]
Da Silva's highlighted exceptional character is
\[
 \alpha=(7,10,13,19,22,28)\pmod {33}.
\]
It is not coordinatewise divisible by $11$, so Theorem~\ref{thm:sector}
independently forces its $W_{11}$-projection to vanish.  Corollary
\ref{cor:class} is stronger: every primitive projection of $W_{11}$ vanishes.

Thus Question~1 of \cite{daSilva} has a negative answer for the proposed
surface $W$.  Nothing here proves the exceptional character nonalgebraic, nor
does it settle the Hodge conjecture for the degree-$33$ Fermat fourfold or in
general.

The version of record of \cite{daSilva} has the correct term $3e_3$ in
Newton's identity.  An arXiv version displays $3e_3^3$, which is
dimensionally inhomogeneous and is evidently a typographical error.  Our
argument uses the corrected published formula.

\section*{Reproducibility and disclosure}
The accompanying dependency-free verifier checks the degree-$33$ character,
the numerical cycle-class specialization, and the symmetry obstruction using
exact integer arithmetic.  The proof itself is symbolic and does not depend on
the program.

This result was found, drafted, and internally audited by an autonomous AI
mathematics system.  Multiple independent AI agents reconstructed the ideal
identity, regular-sequence argument, Chow computation, symmetry obstruction,
and source wording; hostile audits specifically tested scheme structure,
embedded components, character conventions, and scope.  No human endorsement
is claimed.

\begin{thebibliography}{9}
\bibitem{daSilva}
G. da Silva Jr., \emph{Notes on the Hodge conjecture for Fermat varieties},
Experimental Results \textbf{2} (2021), e22,
\href{https://doi.org/10.1017/exp.2021.14}{doi:10.1017/exp.2021.14}.
\bibitem{Fulton}
W. Fulton, \emph{Intersection Theory}, 2nd ed., Springer, 1998.
\bibitem{Shioda}
T. Shioda, \emph{The Hodge conjecture for Fermat varieties}, Math. Ann.
\textbf{245} (1979), 175--184.
\end{thebibliography}
\end{document}
